% =====================================================================
% 02_setup.tex — households, optimality, market clearing, equilibrium (WP-A)
% =====================================================================
\section{Setup: households, optimality, market clearing, equilibrium}
\label{sec:setup}

This section transcribes the building blocks of the model needed for the proofs:
the Home household problem, the convenience-yield wedge, the supply side and
policy, the household first-order conditions and pricing kernel, the auxiliary
variable definitions, the re-scaled (stationary) recursive equilibrium, and the
Devereux--Sutherland portfolio method. We follow KL's equation numbering
throughout (``their eq.'' for the main paper, ``their app.\ eq.'' for the online
appendix) so that the proof work packages can cross-reference. This is setup,
not yet a proof: where an object is stated rather than derived, that is by
design.

\subsection{The Home household problem}
\label{sec:hh-problem}

The representative Home household has recursive (Epstein--Zin) preferences over
consumption $c_t$, labor $\ell_t$, and the real value of safe dollar bonds
$B_{Ht,s}/P_t$ (their eq.~(1)):
\begin{equation}
\label{eq:ez-value}
\val_t=\left\{(1-\disc)\Bigl[c_t\,\Lphi(\ell_t)\,
\bondu_t\!\bigl(B_{Ht,s}/P_t\bigr)\Bigr]^{1-1/\ies}
+\disc\,\Et\!\Bigl[\val_{t+1}^{1-\rra}\Bigr]^{\frac{1-1/\ies}{1-\rra}}\right\}^{\frac{1}{1-1/\ies}}.
\end{equation}
Consumption is a CES aggregate of Home- and Foreign-produced goods with
home-bias weight $\hbias$ and trade elasticity $\tel$ (their eq.~(2)):
\begin{equation}
\label{eq:ces-c}
c_t=\left[\Bigl(\tfrac{1}{1+\zs}+\hbias\Bigr)^{\!\frac1\tel}c_{Ht}^{\frac{\tel-1}{\tel}}
+\Bigl(\tfrac{\zs}{1+\zs}-\hbias\Bigr)^{\!\frac1\tel}c_{Ft}^{\frac{\tel-1}{\tel}}\right]^{\frac{\tel}{\tel-1}}.
\end{equation}
The disutility of labor follows Shimer (2010) and Trabandt--Uhlig (2011)
(their eq.~(3)):
\begin{equation}
\label{eq:labor-phi}
\Lphi(\ell_t)=\left[1+(1/\ies-1)\,\labdis\,
\frac{(\ell_t)^{1+1/\frisch}}{1+1/\frisch}\right]^{\frac{1/\ies}{1-1/\ies}}.
\end{equation}
The function $\bondu_t(\cdot)$ delivers the nonpecuniary (liquidity/safety) value
of safe dollar bonds, in the spirit of money-in-the-utility (Sidrauski 1967) and
the convenience-yield literature (Krishnamurthy--Vissing-Jorgensen 2012); its
functional form is given in \eqref{eq:Omega-form} below. The household supplies a
continuum of labor varieties $j\in[0,1]$, $\ell_t=\int_0^1\ell_t(j)\,dj$.

Each period the household chooses consumption and a portfolio
$\{B_{Ht,s},B_{Ht,o},B_{Ft},k_t\}$ --- safe dollar bonds (nominal rate $i_t$),
other dollar bonds (rate $\ioth_t$), Foreign nominal bonds (rate $i_t^{*}$), and
capital (price $Q^k_t$, dividend $\Pi_{t+1}$, depreciation $\delta$) --- subject
to the resource constraint (their eq.~(4)):
\begin{align}
\label{eq:budget}
P_{Ht}c_{Ht}+E_t^{-1}P^{*}_{Ft}c_{Ft}
&+B_{Ht,s}+B_{Ht,o}+E_t^{-1}B_{Ft}+Q^k_t k_t \notag\\
\le\;(1+i_{t-1})B_{Ht-1,s}&+(1+\ioth_{t-1})B_{Ht-1,o}+E_t^{-1}(1+i^{*}_{t-1})B_{Ft-1}\notag\\
&+\Bigl[\Pi_t+(1-\dep)Q^k_t\Bigr]k_{t-1}\exp(\varphi_t)
+\int_0^1 W_t(j)\ell_t(j)\,dj\notag\\
&\quad-\int_0^1 AC^W_t(j)\,dj+T_t,
\end{align}
where $P_{Ht},P^{*}_{Ft}$ are Home/Foreign producer prices in their own units of
account, $E_t$ is the nominal exchange rate (Foreign units per dollar), the
disaster term $\exp(\varphi_t)$ scales the capital stock, $T_t$ is a government
transfer, and producer-currency pricing is assumed so the law of one price
holds. The Rotemberg wage-adjustment cost is (their eq.~(5))
\begin{equation}
\label{eq:rotemberg}
AC^W_t(j)=\frac{\wadj}{2}\,W_t\ell_t
\left[\frac{W_t(j)}{W_{t-1}(j)\exp(\varphi_t)}-1\right]^2,
\end{equation}
rebated lump-sum. In the simplified environment of
Definition~\ref{def:simplified}, \eqref{eq:rotemberg} is replaced by flexible or
one-period-predetermined wages.

\paragraph{Foreign analog.} Foreign households (measure $\zs$) solve an
analogous problem with risk aversion $\rraF$ that may differ from $\rra$, the
same $\ies,\hbias,\tel,\frisch,\wadj$, and possibly a different discount factor
$\discF$ and labor-disutility scale $\bar\nu^{*}$. They too receive nonpecuniary
utility $\bondu^{*}_t\bigl(B^{*}_{Ht,s}/(E_t^{-1}P^{*}_t)\bigr)$ from safe dollar
bonds. All starred objects are the Foreign counterparts of the Home objects
above.

\subsection{The convenience yield}
\label{sec:cy}

Because safe dollar bonds and other dollar bonds pay in the same unit of account
and are both risk-free, investor indifference at Home requires that the
\emph{effective} return on safe bonds equal the return on other bonds
(their eq.~(14)):
\begin{equation}
\label{eq:wedge}
\frac{1+i_t}{\,1-c_t\,\bondu_t'(B_{Ht,s}/P_t)/\bondu_t(B_{Ht,s}/P_t)\,}=1+\ioth_t.
\end{equation}
The denominator's second term is the marginal nonpecuniary value of safe bonds.
Equating this nonpecuniary value across Home and Foreign agents on the margin
defines the (common) convenience yield $\cy_t$ (their eq.~(15)):
\begin{equation}
\label{eq:omega-def}
\cy_t\equiv c_t\,\frac{\bondu_t'(B_{Ht,s}/P_t)}{\bondu_t(B_{Ht,s}/P_t)}
=c_t^{*}\,\frac{\bondu_t^{*}{}'\!\bigl(B^{*}_{Ht,s}/(E_t^{-1}P^{*}_t)\bigr)}
{\bondu_t^{*}\!\bigl(B^{*}_{Ht,s}/(E_t^{-1}P^{*}_t)\bigr)}.
\end{equation}
KL adopt the convenient functional form (their displayed equation on p.~1659):
\begin{equation}
\label{eq:Omega-form}
\bondu_t\!\left(\frac{B_{Ht,s}}{P_t}\right)
=\exp\!\left(\omega^d_t\frac{B_{Ht,s}}{P_t\,\bar c_t}
-\frac12\frac1{\edem}\Bigl(\frac{B_{Ht,s}}{P_t\,\bar c_t}\Bigr)^{2}
-\Bigl[\omega^d_t\frac{\bar B_{Ht,s}}{P_t\,\bar c_t}
-\frac12\frac1{\edem}\Bigl(\frac{\bar B_{Ht,s}}{P_t\,\bar c_t}\Bigr)^{2}\Bigr]\right),
\end{equation}
where $\omega^d_t$ is an exogenous driving force (the \emph{safety shock}),
$\edem$ a parameter, and bars denote aggregates the household takes as given. The
first two terms inside the exponential make the nonpecuniary value
time-varying; the bracketed terms ensure $\bondu_t(B_{Ht,s}/P_t)=1$ in
equilibrium, so the convenience yield has no mechanical effect on the
household's position through the level of $\bondu_t$. Given the analogous Foreign
form, the second equality in \eqref{eq:omega-def} together with market clearing
in safe dollar bonds implies (their eq.~(16))
\begin{equation}
\label{eq:safe-bond-clearing}
\frac{B_{Ht,s}}{P_t c_t}=\frac{B^{*}_{Ht,s}}{E_t^{-1}P^{*}_t c^{*}_t}
=\frac{(-B^g_{Ht,s})}{P_t c_t+\zs E_t^{-1}P^{*}_t c^{*}_t},
\end{equation}
and the first equality in \eqref{eq:omega-def} then yields the
\emph{reduced-form} convenience yield (their eq.~(17))
\begin{equation}
\label{eq:omega-reduced}
\cy_t=\omega^d_t-\frac1{\edem}\,
\frac{(-B^g_{Ht,s})}{P_t c_t+\zs E_t^{-1}P^{*}_t c^{*}_t}.
\end{equation}
Intuitively, $\cy_t$ rises with private demand for safe dollar bonds
$\omega^d_t$ (latent demand) and falls with public supply $-B^g_{Ht,s}$; the
relative strength of supply is governed by $\edem$, the elasticity of demand to
the nonpecuniary value. Under the fiscal rule \eqref{eq:fiscal-supply} below the
supply term is proportional to global consumption, so $\cy_t$ inherits the
stochastic properties of the safety shock $\omega^d_t$ and is effectively
exogenous.

\subsection{Supply side and policy}
\label{sec:supply-policy}

\paragraph{Producers (their app.\ A.5.3).} The representative Home producer hires
$\ell_t$ from the domestic labor packer and rents $\kappa_t$ of capital on the
international market, producing with productivity $z_t$ and capital share
$\alpha$ (their eq.~(7)). Its first-order conditions are
\begin{equation}
\label{eq:prod-foc-home}
\rw_t=\frac{P_{Ht}}{P_t}(1-\alpha)z_t^{1-\alpha}\ell_t^{-\alpha}\kappa_t^{\alpha},
\qquad
\divr_t=\frac{P_{Ht}}{P_t}\alpha z_t^{1-\alpha}\ell_t^{1-\alpha}\kappa_t^{\alpha-1}.
\end{equation}
The representative Foreign producer's FOCs are analogous, with relative
productivity $z_{Ft}$ and the real exchange rate $q_t$ converting units:
\begin{equation}
\label{eq:prod-foc-foreign}
\rw_t^{*}=q_t^{-1}\frac{P^{*}_{Ft}}{P^{*}_t}(1-\alpha)(z_t z_{Ft})^{1-\alpha}
(\zs\ell_t^{*})^{-\alpha}\kappa_t^{*\alpha},\qquad
\divr_t=q_t^{-1}\frac{P^{*}_{Ft}}{P^{*}_t}\alpha(z_t z_{Ft})^{1-\alpha}
(\zs\ell_t^{*})^{1-\alpha}\kappa_t^{*\alpha-1}.
\end{equation}
Because capital is rented on a single international market, its dividend
$\divr_t$ (in dollars) is common across countries.

\paragraph{Unions (their app.\ A.5.2).} Each Home union $j$ sets the wage
$W_t(j)$ and supplies labor $\ell_t(j)$ to maximize members' utilitarian
welfare subject to the Rotemberg cost \eqref{eq:rotemberg}. The resulting wage
first-order condition is
\begin{equation}
\label{eq:union-foc}
\begin{aligned}
\rw_t-\frac{c_t\Lphi'(\ell_t)}{\Lphi(\ell_t)}
+\rw_t\frac{\wadj}{\lel}\Bigg[
&\frac{\rw_t}{\rw_{t-1}\exp(\varphi_t)}\frac{P_t}{P_{t-1}}
\left(\frac{\rw_t}{\rw_{t-1}\exp(\varphi_t)}\frac{P_t}{P_{t-1}}-1\right)\\
&-\Et m_{t,t+1}\left(\frac{\rw_{t+1}}{\rw_t\exp(\varphi_{t+1})}\right)^{2}
\frac{P_{t+1}}{P_t}\frac{\ell_{t+1}}{\ell_t}
\left(\frac{\rw_{t+1}}{\rw_t\exp(\varphi_{t+1})}\frac{P_{t+1}}{P_t}-1\right)
\Bigg]=0,
\end{aligned}
\end{equation}
with an analogous condition in Foreign. In the simplified environment this
collapses to a flexible or one-period-predetermined wage equation.

\paragraph{Global capital producer (their eqs.~(8)--(9)).} A representative
global capital producer combines Home and Foreign consumption goods (without home
bias) to produce new capital with adjustment-cost scale $\chi^x$; profit
maximization yields the real price of capital
\begin{equation}
\label{eq:qk}
\qk_t=\left(\frac{\bar k_t}{\bar k_{t-1}\exp(\varphi_t)}\right)^{\chi^x}
\left[\Bigl(\tfrac{1}{1+\zs}+\hbias\Bigr)\Bigl(\tfrac{P_{Ht}}{P_t}\Bigr)^{1-\tel}
+\Bigl(\tfrac{\zs}{1+\zs}-\hbias\Bigr)\Bigl(\tfrac{P^{*}_{Ft}}{P_t}\Bigr)^{1-\tel}\right]^{\frac1{1-\tel}}.
\end{equation}
In the simplified environment $\chi^x\!\to\!\infty$ with $\bar k_t$ fixed.

\paragraph{Monetary policy (their eqs.~(10)--(11)).} A Taylor (1993) rule
governs the Home nominal rate,
\begin{equation}
\label{eq:taylor}
1+i_t=(1+\bar\imath)\left(\frac{P_t}{P_{t-1}}\right)^{\taylor},
\end{equation}
where the ideal CPI is
\begin{equation}
\label{eq:cpi}
P_t=\left[\Bigl(\tfrac{1}{1+\zs}+\hbias\Bigr)P_{Ht}^{1-\tel}
+\Bigl(\tfrac{\zs}{1+\zs}-\hbias\Bigr)\bigl(E_t^{-1}P^{*}_{Ft}\bigr)^{1-\tel}\right]^{\frac1{1-\tel}}.
\end{equation}
An analogous CPI-targeting rule with the same coefficient $\taylor$ determines
$i_t^{*}$ in Foreign.

\paragraph{Fiscal policy (their eqs.~(12)--(13)).} The Home government maintains
a constant ratio of safe dollar debt to global consumption,
\begin{equation}
\label{eq:fiscal-supply}
-B^g_{Ht,s}=\bar b^g\bigl(P_t c_t+\zs E_t^{-1}P^{*}_t c^{*}_t\bigr),
\end{equation}
with $\bar b^g>0$ the empirically relevant case (the government borrows in safe
dollar bonds, i.e.\ issues Treasuries), and rebates with lump-sum transfers
\begin{equation}
\label{eq:fiscal-transfer}
T_t=\int_0^1 AC^W_t(j)\,dj+(1+i_{t-1})B^g_{Ht-1,s}-B^g_{Ht,s}.
\end{equation}
Ricardian equivalence holds for all government activity other than safe-bond
issuance, which alone provides nonpecuniary services. Combining
\eqref{eq:fiscal-supply} with \eqref{eq:omega-reduced} renders $\cy_t$
proportional to the exogenous safety shock $\omega^d_t$ (up to the
constant $\bar b^g/\edem$).

\subsection{Household first-order conditions and pricing kernel}
\label{sec:hh-foc}

Home intratemporal optimality across Home and Foreign goods is (their app.\
A.5.1)
\begin{equation}
\label{eq:intratemporal}
\frac{c_{Ht}}{c_{Ft}}
=\frac{\tfrac{1}{1+\zs}+\hbias}{\tfrac{\zs}{1+\zs}-\hbias}\,\tot_t^{-\tel},
\qquad \tot_t\equiv\frac{E_t P_{Ht}}{P^{*}_{Ft}}\ \ \text{(terms of trade)}.
\end{equation}
The Home \emph{real pricing kernel} is the Epstein--Zin stochastic discount
factor (their app.\ A.5.1), evaluated using $\bondu_t(B_{Ht,s}/P_t)=1$ in
equilibrium:
\begin{equation}
\label{eq:kernel}
\pk_{t,t+1}=\disc\,\frac{c_t}{c_{t+1}}
\left(\frac{c_{t+1}\Lphi(\ell_{t+1})}{c_t\Lphi(\ell_t)}\right)^{1-1/\ies}
\left(\frac{\val_{t+1}}{\ce_t}\right)^{1/\ies-\rra},
\qquad
\ce_t=\Et\!\Bigl[\val_{t+1}^{1-\rra}\Bigr]^{\frac1{1-\rra}},
\end{equation}
where $\ce_t$ is the certainty equivalent of the continuation value. Intertemporal
optimality is characterized by three Euler equations (their app.\ A.5.1) --- for
safe dollar bonds (adjusted by the convenience yield), Foreign bonds (UIP), and
capital, respectively:
\begin{align}
\label{eq:euler-safe}
1&=\Et\,\pk_{t,t+1}\,\frac{1+\rr_{t+1}}{1-\cy_t},\\
\label{eq:euler-uip}
1&=\Et\,\pk_{t,t+1}\,\frac{\rer_t}{\rer_{t+1}}\bigl(1+\fr{r}_{t+1}\bigr),\\
\label{eq:euler-capital}
1&=\Et\,\pk_{t,t+1}\bigl(1+\rk_{t+1}\bigr).
\end{align}
Equation \eqref{eq:euler-safe} is implied by optimality in \emph{either} safe or
other dollar bonds given the definition of $\cy_t$ in \eqref{eq:omega-def}: the
$1/(1-\cy_t)$ factor converts the safe-bond return into the convenience-adjusted
return that the kernel prices. Analogous starred conditions hold for Foreign
households, using the Foreign kernel $\pk^{*}_{t,t+1}$ (their app.\ eqs.\
(47)--(49)).

\paragraph{Safe-bond market clearing and the real resource constraint.} Combining
the cross-country indifference condition \eqref{eq:omega-def} (which gives
$B_{Ht,s}/(P_t c_t)=B^{*}_{Ht,s}/(E_t^{-1}P^{*}_t c^{*}_t)$) with global market
clearing in safe dollar bonds yields the allocation of safe bonds across
countries (their app.\ A.5.1, restating \eqref{eq:safe-bond-clearing}):
\begin{equation}
\label{eq:safe-bond-split}
B_{Ht,s}=\frac{P_t c_t}{P_t c_t+\zs E_t^{-1}P^{*}_t c^{*}_t}(-B^g_{Ht,s}),
\qquad
B^{*}_{Ht,s}=\frac{E_t^{-1}P^{*}_t c^{*}_t}{P_t c_t+\zs E_t^{-1}P^{*}_t c^{*}_t}(-B^g_{Ht,s}).
\end{equation}
Substituting government transfers \eqref{eq:fiscal-transfer} into the budget
constraint \eqref{eq:budget}, dividing by $P_t$, and denoting real positions
$b_{Ht,s}\equiv B_{Ht,s}/P_t$, $b^g_{Ht,s}\equiv B^g_{Ht,s}/P_t$,
$b_{Ht,o}\equiv B_{Ht,o}/P_t$, $b_{Ft}\equiv B_{Ft}/P_t$, $\qk_t\equiv Q^k_t/P_t$,
$\divr_t\equiv\Pi_t/P_t$, $\rw_t\equiv W_t/P_t$, the Home household's real
resource constraint is (their app.\ A.5.1)
\begin{equation}
\label{eq:real-budget-home}
\begin{aligned}
c_t+b_{Ht,s}+b^g_{Ht,s}+b_{Ht,o}+\rer_t^{-1}b_{Ft}+\qk_t k_t
&=\rw_t\ell_t
+(1+\rr_t)\bigl(b_{Ht-1,s}+b^g_{Ht-1,s}\bigr)\\
&\quad+\left(\frac{1+\rr_t}{1-\cy_{t-1}}\right)b_{Ht-1,o}
+\rer_t^{-1}(1+\fr{r}_t)b_{Ft-1}\\
&\quad+\bigl(\divr_t+(1-\dep)\qk_t\bigr)k_{t-1}\exp(\varphi_t).
\end{aligned}
\end{equation}
KL further define each agent's net dollar-bond position
$b_{Ht}\equiv(1-\cy_t)(b_{Ht,s}+b^g_{Ht,s})+b_{Ht,o}$ and
$b^{*}_{Ht}\equiv(1-\cy_t)b^{*}_{Ht,s}+b^{*}_{Ht,o}$, each earning the
``other-bond'' return $1+\ioth_t=(1+i_t)/(1-\cy_t)$; the composition of the
dollar-bond position matters only for the seigniorage Home earns on the safe debt
held by Foreign (their p.~1661 and app.\ A.5.1).

\subsection{Additional variables (their app.\ A.4)}
\label{sec:aux-vars}

The following auxiliary definitions are used throughout the proofs. The real
exchange rate (an increase is a Home appreciation):
\begin{equation}
\label{eq:rer-def}
\rer_t\equiv\frac{E_t P_t}{P^{*}_t}.
\end{equation}
The real interest rates and the real return on capital (in Home consumption
units):
\begin{equation}
\label{eq:returns-def}
1+\rr_{t+1}\equiv(1+i_t)\frac{P_t}{P_{t+1}},\quad
1+\fr{r}_{t+1}\equiv(1+i^{*}_t)\frac{P^{*}_t}{P^{*}_{t+1}},\quad
1+\rk_{t+1}\equiv\frac{\bigl(\Pi_{t+1}+(1-\dep)Q^k_{t+1}\bigr)}{Q^k_t}\frac{P_t}{P_{t+1}}\exp(\varphi_{t+1}).
\end{equation}
Output, the real value of aggregate saving, and real net foreign assets at Home:
\begin{equation}
\label{eq:agg-def}
\out_t\equiv(z_t\ell_t)^{1-\alpha}(\kappa_t)^{\alpha},\quad
\sav_t\equiv\frac{1}{P_t}\bigl(B_{Ht}+E_t^{-1}B_{Ft}+Q^k_t k_t\bigr),\quad
\nfa_t\equiv\sav_t-\frac{Q^k_t}{P_t}\kappa_{t+1}\exp(-\varphi_{t+1}).
\end{equation}
All variables are dated at the end of the period (their app.\ fn.~53); the factor
$\exp(-\varphi_{t+1})$ in $\nfa_t$ undoes capital destruction so that capital
\emph{used} at Home is compared with capital \emph{owned} by Home residents (their
app.\ fn.~54).

\subsection{The re-scaled (stationary) recursive equilibrium}
\label{sec:rescaled}

Because global productivity $z_t$ has a unit root (their eq.~(18)), KL stationarize
by dividing all real variables (and real-deflated nominal variables) by $z_t$,
\emph{labor excepted} (their app.\ A.6). Write $\rs{x}\equiv x/z_t$; lagged stocks
carry a second deflation by the growth term
$\exp(\sigma^z\epsilon^z_t+\varphi_t)$, denoted with a double tilde
$\widetilde{\widetilde{x}}_{t-1}$. After scaling, global-productivity shocks
(inclusive of disasters) only govern the transition across states. The recursive
equilibrium is the system of their app.\ eqs.~(33)--(72); we transcribe the
economically central conditions and refer to the rest by number.

\paragraph{Home block (their app.\ eqs.~(33)--(41)).} The re-scaled value
recursion, certainty equivalent, CES aggregator, intratemporal allocation, and
kernel are their app.\ eqs.~(33)--(37). The three re-scaled Euler equations are
\begin{align}
\label{eq:reuler-safe}
1&=\Et\,\rs{\pk}_{t,t+1}\exp\!\bigl(-\rra[\sigma^z\epsilon^z_{t+1}+\varphi_{t+1}]\bigr)
\frac{1+\rr_{t+1}}{1-\cy_t}, &&\text{(their app.\ eq.~(38))}\\
\label{eq:reuler-uip}
1&=\Et\,\rs{\pk}_{t,t+1}\exp\!\bigl(-\rra[\sigma^z\epsilon^z_{t+1}+\varphi_{t+1}]\bigr)
\frac{\rer_t}{\rer_{t+1}}(1+\fr{r}_{t+1}), &&\text{(their app.\ eq.~(39))}\\
\label{eq:reuler-cap}
1&=\Et\,\rs{\pk}_{t,t+1}\exp\!\bigl(-\rra[\sigma^z\epsilon^z_{t+1}+\varphi_{t+1}]\bigr)
(1+\rk_{t+1}), &&\text{(their app.\ eq.~(40))}
\end{align}
and the re-scaled Home resource constraint is
\begin{equation}
\label{eq:rbudget-home}
\rs{c}_t+\rs{b}_{Ht}+\rer_t^{-1}\rs{b}_{Ft}+\qk_t\rs{k}_t
=\rs{w}_t\ell_t+\theta_t\bigl(\divr_t+(1-\dep)\qk_t\bigr)\widetilde{\widetilde{k}}_{t-1},
\quad\text{(their app.\ eq.~(41))}
\end{equation}
where $\theta_t$ is the Home wealth share. The Foreign block (their app.\
eqs.~(42)--(50)) is symmetric, with risk aversion $\rraF$, Euler equations
\begin{align}
\label{eq:reuler-safe-f}
1&=\Et\,\rs{\pk}^{*}_{t,t+1}\exp\!\bigl(-\rraF[\sigma^z\epsilon^z_{t+1}+\varphi_{t+1}]\bigr)
\frac{\rer_{t+1}}{\rer_t}\frac{1+\rr_{t+1}}{1+\cy_t}, &&\text{(their app.\ eq.~(47))}\\
\label{eq:reuler-uip-f}
1&=\Et\,\rs{\pk}^{*}_{t,t+1}\exp\!\bigl(-\rraF[\sigma^z\epsilon^z_{t+1}+\varphi_{t+1}]\bigr)
(1+\fr{r}_{t+1}), &&\text{(their app.\ eq.~(48))}\\
\label{eq:reuler-cap-f}
1&=\Et\,\rs{\pk}^{*}_{t,t+1}\exp\!\bigl(-\rraF[\sigma^z\epsilon^z_{t+1}+\varphi_{t+1}]\bigr)
\frac{\rer_{t+1}}{\rer_t}(1+\rk_{t+1}), &&\text{(their app.\ eq.~(49))}
\end{align}
and the Foreign resource constraint (their app.\ eq.~(50)).

\noindent\emph{(Note: these match KL's app.\ eqs.\ (47)--(49) verbatim, verified
against their appendix p.~64. Observe the Home/Foreign asymmetry in the
convenience adjustment: the Home safe-bond Euler equation~\eqref{eq:reuler-safe}
carries $1/(1-\cy_t)$, whereas the Foreign safe-bond Euler
equation~\eqref{eq:reuler-safe-f} carries $1/(1+\cy_t)$, because the convenience
benefit accrues to the issuer's residents and is netted out of the price Foreign
agents are willing to pay. The $q_{t+1}/q_t$ placement is exactly as in KL.)}

\paragraph{Wealth-share law of motion (their app.\ eq.~(51)).} The global wealth
share of Home households, inclusive of seigniorage, evolves as
\begin{equation}
\label{eq:wealth-share}
\theta_{t+1}=\frac{1}{(\divr_{t+1}+(1-\dep)\qk_{t+1})\widetilde{\widetilde{k}}_t}
\left[\frac{1+\rr_{t+1}}{1-\cy_t}\rs{b}_{Ht}
+\frac{1}{\rer_{t+1}}(1+\fr{r}_{t+1})\rs{b}_{Ft}
+(\divr_{t+1}+(1-\dep)\qk_{t+1})\widetilde{\widetilde{k}}_t
-\cy_{t+1}\frac{\zs\rer^{-1}_{t+1}\rs{c}^{*}_{t+1}}{\rs{c}_{t+1}+\zs\rer^{-1}_{t+1}\rs{c}^{*}_{t+1}}\rs{b}^g_{Ht+1,s}\right].
\end{equation}
The last term is the seigniorage Home earns on the share of its safe debt held by
Foreign.

\paragraph{Supply, market clearing, returns, prices (their app.\ eqs.~(52)--(72)).}
Supply-side optimality (re-scaled union and producer FOCs) is their app.\
eqs.~(52)--(60); goods-, factor-, and asset-market clearing
\begin{align}
\label{eq:mc-goods}
&\rs{c}_{Ht}+\zs\rs{c}^{*}_{Ht}+\Bigl(\tfrac{\widetilde{\widetilde{k}}_t}{\widetilde{\widetilde{k}}_{t-1}}\Bigr)^{\!\chi^x}\rs{x}_{Ht}=(\ell_t)^{1-\alpha}(\rs\kappa_t)^{\alpha}, &&\text{(their app.\ eq.~(61))}\\
&\rs{c}_{Ft}+\zs\rs{c}^{*}_{Ft}+\Bigl(\tfrac{\widetilde{\widetilde{k}}_t}{\widetilde{\widetilde{k}}_{t-1}}\Bigr)^{\!\chi^x}\rs{x}_{Ft}=(z_{Ft}\zs\ell^{*}_t)^{1-\alpha}(\rs\kappa^{*}_t)^{\alpha}, &&\text{(their app.\ eq.~(62))}\\
&\rs\kappa_t+\rs\kappa^{*}_t=\widetilde{\widetilde{k}}_{t-1}, \quad
\rs{k}_t+\zs\rs{k}^{*}_t=\widetilde{\widetilde{k}}_t,\quad
(1-\dep)\widetilde{\widetilde{k}}_{t-1}+\rs{x}_t=\widetilde{\widetilde{k}}_t, &&\text{(their app.\ eqs.~(63)--(65))}\\
\label{eq:mc-bonds}
&\rs{b}_{Ht}+\zs\rs{b}^{*}_{Ht}=0, &&\text{(their app.\ eq.~(66))}
\end{align}
the return definitions
\begin{equation}
\label{eq:rescaled-returns}
1+\rr_{t+1}=(1+i_t)\tfrac{P_t}{P_{t+1}},\quad
1+\fr{r}_{t+1}=(1+i^{*}_t)\tfrac{P^{*}_t}{P^{*}_{t+1}},\quad
1+\rk_{t+1}=\tfrac{\divr_{t+1}+(1-\dep)\qk_{t+1}}{\qk_t}\exp(\varphi_{t+1}),
\end{equation}
(their app.\ eqs.~(67)--(69)), and the price relations linking the CPI, terms of
trade, and real exchange rate (their app.\ eqs.~(70)--(72)):
\begin{equation}
\label{eq:price-relations}
\frac{P_t}{P_{Ht}}=\Bigl[\bigl(\tfrac{1}{1+\zs}+\hbias\bigr)+\bigl(\tfrac{\zs}{1+\zs}-\hbias\bigr)\tot_t^{\tel-1}\Bigr]^{\frac1{1-\tel}},
\quad
\rer_t=\tot_t\left[\frac{(\tfrac{1}{1+\zs}+\hbias)+(\tfrac{\zs}{1+\zs}-\hbias)\tot_t^{\tel-1}}{(\tfrac{1}{1+\zs}-\tfrac{\hbias}{\zs})\tot_t^{1-\tel}+(\tfrac{\zs}{1+\zs}+\tfrac{\hbias}{\zs})}\right]^{\frac1{1-\tel}}.
\end{equation}

\paragraph{Equilibrium and state vector.} Together with the Taylor and fiscal
rules and the driving forces, their app.\ eqs.~(33)--(72) define the recursive
equilibrium in the sense of the following definition.

\begin{definition}[Equilibrium; KL app.\ A.3 Definition~1]
\label{def:equilibrium}
An equilibrium is a sequence of prices and policies such that: each Home
representative household chooses $\{c_{Ht},c_{Ft},B_{Ht,s},B_{Ht,o},B_{Ft},k_t\}$
to maximize \eqref{eq:ez-value} subject to \eqref{eq:ces-c}--\eqref{eq:rotemberg}
(and analogously in Foreign); each Home union $j$ chooses $\{W_t(j),\ell_t(j)\}$
to maximize members' utilitarian welfare subject to \eqref{eq:rotemberg}
(analogously in Foreign); the Home labor packer chooses $\{\ell_t(j)\}$ to
maximize profits (their eq.~(6), analogously in Foreign); the Home producer
chooses $\{\ell_t,\kappa_t\}$ to maximize profits (their eq.~(7), analogously in
Foreign); the global capital producer chooses $\{x_{Ht},x_{Ft},x_t\}$ to maximize
profits (their eq.~(9)) subject to \eqref{eq:qk}; the Home government sets
$B^g_{Ht,s}$ by \eqref{eq:fiscal-supply} and $\{i_t,\{T_t\}\}$ by
\eqref{eq:taylor} and \eqref{eq:fiscal-transfer} (and the Foreign government the
latter); and goods, factor, and asset markets clear (their app.\ eqs.~(23)--(32),
including $B_{Ft}+\zs B^{*}_{Ft}=0$, their app.\ eq.~(32)).
\end{definition}

After re-scaling, the aggregate state in period $t$ is seven-dimensional:
\begin{equation}
\label{eq:state}
\bigl\{\,p_t,\ \cy_t,\ z_{Ft},\ \theta_t,\ \widetilde{\widetilde{k}}_{t-1},\
\rs{w}_{t-1},\ \rs{w}^{*}_{t-1}\,\bigr\},
\end{equation}
i.e.\ disaster probability, convenience yield, Foreign relative productivity,
Home wealth share, lagged (re-scaled) global capital, and lagged (re-scaled) Home
and Foreign real wages. By Walras' law the Foreign bond market clears given the
rest. In the simplified environment, $p$ and $z_F$ are constant
(Definition~\ref{def:simplified}(iv)) and the wage states are trivial under
flexible or one-period-predetermined wages, so the effective state collapses to
$\{\cy,\theta\}$ around the steady state --- the two variables that carry the
analytical dynamics.

\subsection{The Devereux--Sutherland portfolio method}
\label{sec:ds-method}

Equilibrium \emph{portfolios} --- the steady-state holdings of dollar bonds,
capital, and Foreign bonds that each country chooses --- cannot be read off the
first-order (linearized) equilibrium conditions alone. KL pin them down with the
second-order method of Devereux and Sutherland (2011); we lay out the logic here
and defer the algebra to the portfolio work package (WP-D).

To first order, the three asset returns
$\{1+\rr_{t+1},\,1+\rk_{t+1},\,(\rer_t/\rer_{t+1})(1+\fr{r}_{t+1})\}$ enter the
Euler equations \eqref{eq:euler-safe}--\eqref{eq:euler-capital} only through
their conditional means, which the Euler equations equate. Hence to first order
all assets are \emph{perfect substitutes} and the portfolio weights are
\emph{indeterminate}: any portfolio is consistent with the first-order
equilibrium. The Devereux--Sutherland insight is that portfolios are pinned down
by a \emph{second-order} object. Specifically, the steady-state portfolio is the
unique one for which the first-order \emph{excess-return} (risk-sharing)
conditions --- the differences of the Euler equations
\eqref{eq:euler-safe}--\eqref{eq:euler-capital} across assets --- hold once the
portfolio Euler equations are approximated to \emph{second} order. Equivalently,
the portfolio is chosen so that the conditional covariances between the pricing
kernel and excess returns (which are second-order objects) are consistent with
the first-order dynamics of consumption and the real exchange rate that those
same portfolios induce.

Two structural facts make this tractable and economically sharp in the simplified
environment. First, there are exactly \emph{three} traded assets (safe dollar
bonds, capital, Foreign bonds). Second, by Definition~\ref{def:simplified}(iv)
the economy is subject to exactly \emph{two} aggregate shocks (global
productivity $\epsilon^z$ and safety $\epsilon^\omega$). With one more asset than
shock, the asset span is \emph{locally complete} in the sense of Coeurdacier and
Gourinchas (2016): around the deterministic steady state the available assets
suffice to implement the constrained-efficient (complete-markets) risk-sharing
allocation. Consequently the equilibrium portfolios that the
Devereux--Sutherland method selects are precisely those that implement efficient
risk sharing to first order. This is the formal content behind KL's statement
(their p.~1664) that ``the three available assets implement efficient risk
sharing around the steady state.'' The closed-form characterization of these
portfolios and of the resulting risk premia (their Propositions~4--5) is carried
out in WP-D; here we only fix the method and its scope.

\gapflag{KL state local completeness and efficient risk sharing as a consequence
of the ``one more asset than shock'' count (their p.~1664, citing Coeurdacier and
Gourinchas 2016) but do not, in the main text, prove that the
Devereux--Sutherland-selected portfolio coincides with the efficient one in this
specific three-asset/two-shock setting. We state the equivalence here as KL
assert it; the explicit verification (that the span $\{$bond, capital, FX$\}$
covers the two shocks and that no redundancy or rank-deficiency arises at the
symmetric steady state) belongs to WP-D and is flagged for the critic.}
